\section*{Chapter \ref{ch:nw:disaster-recovery-and-regulatory-requirements} --- Disaster Recovery and Regulatory Requirements}

\begin{solution}{exo:nw:disaster-recovery-and-regulatory-requirements:1}
Next business day resumption of trading and two-hour resumption of critical SCI systems following a wide-scale disruption, with backup
capabilities sufficiently resilient and geographically diverse.
\end{solution}

\begin{solution}{exo:nw:disaster-recovery-and-regulatory-requirements:2}
The time objective bounds how long a system may be down; the point objective bounds how much recent data may be lost. The first is set
by the runbook, the second by replication.
\end{solution}

\begin{solution}{exo:nw:disaster-recovery-and-regulatory-requirements:3}
It is \qty{34.2}{\kilo\metre} from Secaucus, inside any regional disaster that takes the primary: a storm, a flood or a power failure
would likely take both.
\end{solution}

\begin{solution}{exo:nw:disaster-recovery-and-regulatory-requirements:4}
Replication sets only the recovery point; the warm design's time is its runbook, dominated by starting and restoring servers and loading
data, which faster replication does not change.
\end{solution}

\begin{solution}{exo:nw:disaster-recovery-and-regulatory-requirements:5}
Alternative arrangements to manage outstanding orders and positions, and arrangements to shut down algorithms or systems without creating
disorderly trading conditions.
\end{solution}

\begin{solution}{exo:nw:disaster-recovery-and-regulatory-requirements:6}
Yes on 28 September 2026; no on 1 November 2026, more than twelve months after the test.
\end{solution}

\begin{solution}{exo:nw:disaster-recovery-and-regulatory-requirements:7}
From 20.3\% to 44.1\%: the restore step is the largest, but the other five steps still take about two hours at their medians together.
\end{solution}

\begin{solution}{exo:nw:disaster-recovery-and-regulatory-requirements:8}
A test is not a decision: before Sandy the industry's test had succeeded two days earlier and the markets still did not switch. The firm
needs a runbook with a decision owner and criteria, rehearsed switching, not only connectivity, and the exchanges' own plans.
\end{solution}

\begin{solution}{pb:nw:disaster-recovery-and-regulatory-requirements:1}
\textbf{Part I.}
\begin{enumerate}
\item NY4 \qty{0.3}{\kilo\metre}, Carteret 26.0, Mahwah 34.2, Aurora 1\,191.1; only Aurora is outside.
\item \qty{17.4}{\milli\second} of round trip to Aurora at a \omterm{def:nw:the-north-american-map:factor}{route factor} of 1.5, \qty{0.5}{\milli\second} to Mahwah.
\item To be outside the disasters that could take the primary, as the rules require geographic diversity.
\item \qty{56}{\milli\second}: \qty{50}{\milli\second} of lag plus the one-way delay.
\end{enumerate}
\textbf{Part II.}
\begin{enumerate}[resume]
\item 99.2\%, 64 and 96 minutes.
\item 20.3\%, 148 and 237 minutes.
\item Up to fifteen minutes, the snapshot interval.
\item Starting and restoring servers (45 minutes at the median) and loading data (30).
\end{enumerate}
\textbf{Part III.}
\begin{enumerate}[resume]
\item Plans reasonably designed to resume critical systems within two hours and trading by the next business day, geographically diverse;
  designated members testing with the exchange at least every twelve months.
\item Documented business continuity arrangements per venue, relocation to a back-up site where appropriate, orderly shutdown, management of
  outstanding orders and positions, annual review and test.
\item Firms had to be able to remain within the \omterm{def:nw:disaster-recovery-and-regulatory-requirements:resilience}{impact tolerances} of their \omterm{def:nw:disaster-recovery-and-regulatory-requirements:resilience}{important business services}.
\item That firms and exchanges might not switch even when the backups worked: the decision and the confidence matter as much as the
  equipment.
\end{enumerate}
\textbf{Part IV.}
\begin{enumerate}[resume]
\item \emph{Named result}: with the chapter's runbooks, the hot design resumes within two hours in 99.2\% of recoveries and the warm design
  in 20.3\%; at 70\% and 25\% of chapter~27's USD 3.19 million a year, the difference costs USD 1.44 million a year.
\item Published: the rules, the exchanges' recovery sites, test dates and site coordinates. Assumed: every step's duration, the costs, the
  hazard radius, the replication lag.
\item If the firm's \omterm{def:nw:disaster-recovery-and-regulatory-requirements:resilience}{impact tolerance} for trading is two hours, and a lost day costs more than USD 1.44 million often enough; otherwise a warm
  site with an orderly shutdown may be the proportionate choice.
\item Pre-staged servers, automated restores, parallel steps, and measured, rehearsed runbooks.
\item Named services and tolerances, the runbook, test results with measured times, and records of participation in exchange tests.
\item As a designated member where asked, and otherwise voluntarily, with the same runbook it would use for real.
\item In the resilience rows: the recovery site, its lines and \omterm{def:nw:colocation-products-and-how-they-are-sold:mmr}{cross-connects}, the replication, the runbook and the test calendar.
\item Hours of trading back after a disaster, at a yearly price.
\end{enumerate}
\end{solution}

\begin{solution}{iq:nw:disaster-recovery-and-regulatory-requirements:1}
The \omterm{def:nw:disaster-recovery-and-regulatory-requirements:objectives}{recovery time objective} is how long a system may be down; the \omterm{def:nw:disaster-recovery-and-regulatory-requirements:objectives}{recovery point objective} how much data may be lost. The runbook sets the
first; replication sets the second.
\iqlookfor{Definitions; runbook against replication.}
\end{solution}

\begin{solution}{iq:nw:disaster-recovery-and-regulatory-requirements:2}
Outside the New York metropolitan area's shared hazards, near the venues' own recovery sites and with good lines to them: Aurora is the
natural choice for U.S. futures and BrokerTec; the costs are asynchronous replication and a recovery point.
\iqlookfor{Hazards; the venues' recovery sites; replication trade-off.}
\end{solution}

\begin{solution}{iq:nw:disaster-recovery-and-regulatory-requirements:3}
One authority decides which site is active (a person with criteria, or a quorum), the old primary is fenced before the new one trades, and
sessions use the exchange's separate recovery addresses so both sites cannot trade at once.
\iqlookfor{Single decision; fencing; separate addresses.}
\end{solution}

\begin{solution}{iq:nw:disaster-recovery-and-regulatory-requirements:4}
Cancel-on-disconnect and the exchanges' kill switches stop resting orders; positions are known from drop copies and the clearing firm;
hedges go through the recovery site or a broker by phone; then the runbook brings trading back.
\iqlookfor{Exposure first; drop copies; fallbacks; runbook.}
\end{solution}

\begin{solution}{iq:nw:disaster-recovery-and-regulatory-requirements:5}
From the \omterm{def:nw:disaster-recovery-and-regulatory-requirements:resilience}{impact tolerance}: simulate or measure each design's recovery time against it, price each, and compare the difference with the
expected cost of the extra downtime.
\iqlookfor{Tolerance first; measured times; cost against expected loss.}
\end{solution}

\begin{solution}{iq:nw:disaster-recovery-and-regulatory-requirements:6}
\omterm{def:nw:disaster-recovery-and-regulatory-requirements:resilience}{Important business services} and tolerances; a recovery site outside the primary's hazards near the venues' recovery sites; runbooks with
measured times; orderly shutdown and management of outstanding orders per venue; annual tests with the exchanges and the industry; records
kept as evidence.
\iqlookfor{Rules mapped to arrangements; testing; evidence.}
\end{solution}
